%%%%%%%%%%%%%%%%%%%%%%%%%%%%%%%%%
%
%       EXERCÍCIO
%
%%%%%%%%%%%%%%%%%%%%%%%%%%%%%%%%%

\ifdefstring{\atividade}{prova}{%
    \renewcommand{\valorquestao}{\ValorQ\ pontos}
}%

\begin{exercicioBanco}[\valorquestao]
    Uma indústria de alimentos afirma que o tempo médio de conservação de um de seus produtos, armazenado sob refrigeração, é de 120 dias. Um laboratório independente deseja testar se, na verdade, esse tempo é inferior ao anunciado. Para isso, selecionou uma amostra aleatória de 25 unidades do produto e encontrou uma média de conservação de 117{,}5 dias. Sabe-se que o desvio padrão populacional é conhecido e igual a 5 dias. Considere que a variável ``tempo de conservação'' segue uma distribuição normal. Com base nessas informações, teste a hipótese unilateral:
%
\[
    H_0: \mu = 120 \quad \text{versus} \quad H_1: \mu < 120
\]
%
ao nível de significância de 1\%. Pergunta: Com base no teste, a hipótese nula deve ser rejeitada ou não?
\end{exercicioBanco}

